\documentclass[10pt]{article}
\usepackage[margin=20mm]{geometry}
\usepackage{amsmath,amssymb,mathtools,fontspec,unicode-math,xurl,hyperref}
\IfFontExistsTF{DejaVu Serif}{\setmainfont{DejaVu Serif}}{\setmainfont{Latin Modern Roman}}
\setmathfont{latinmodern-math.otf}
\setlength{\parindent}{0pt}\setlength{\parskip}{6pt}\sloppy
\title{Technical Review of the Endpoint Obstruction and Complex-Parameter Identifications in a Cusp-Closing Construction}
\author{AI-assisted technical review}
\date{}
\begin{document}
\maketitle
Abstract: The cusp endpoint equation forces a logarithmic term in the inhomogeneous solution used to construct the obstruction space. Consequently Lemma 4.8’s pure-power bound is incompatible with its stated forcing, and its barrier proof loses the relevant power of the endpoint variable. This is a Major defect in a central proof dependency, not a disproof of the main degeneration theorem. A second, local Minor defect concerns a missing complex phase in the expanded identification map; the defining composition and its norm estimates admit an explicit repair. Algebraic smoothness and the principal neck scaling are checked separately. Complete nonlinear gluing and independent acceptance remain unverified.\par
\section*{Source, quantifiers, and scope}
M25-0380 is Xin Fu, Hans-Joachim Hein and Xumin Jiang, A Continuous Cusp Closing Process for Negative Kähler–Einstein Metrics, Geometric and Functional Analysis 35 (2025), 542–632, DOI 10.1007/s00039-025-00708-y [source]. The reviewed source is the 91-page publisher version of record, published 4 March 2025 [publisher]. All page references below use printed journal pagination; physical PDF page \(p\) is printed page \(541+p\). An associated preprint is not substituted for this exact source.\par
Equation reference: eq:family. The family is the projective closure of \[X_\sigma\cap\mathbb C^3
 =\left\{\sum_{j=1}^3(z_j^6+z_j^3)=\sigma\right\},
 \qquad \sigma\in\Delta\subset\mathbb C.
 \] The Main Theorem, p. 544, quantifies over each fixed \(R>0\) and sufficiently small \(0<|\sigma|\) depending on \(R\). After the local analytic coordinate change and scaling by \(\sigma^{-1/3}\), multiplying the metric by \(|\log|\sigma||^{3/2}\) is asserted to give a Tian–Yau limit with \(C^0\) error \(O_{R,\varepsilon}(|\log|\sigma||^{-1+\varepsilon})\) for every \(\varepsilon>0\). The additional claims about the only bubble, curvature loss and topology are distinct assertions. They are not certified by the endpoint calculations here.\par
In the neck analysis, \(n\) is the complex dimension and the application has \(n=2\). The coordinates are \(s=1-t/T\) and \(\eta=1-s=t/T\). Thus \(\eta\downarrow0\) is the cusp endpoint, while \(s\downarrow0\) is the Tian–Yau endpoint. The limiting radial operator is denoted by \(L_\infty\). The review concentrates on its homogeneous and inhomogeneous equations, the exact complex-parameter maps, and the dependencies that carry these into Sections 4–5. The complete source was parsed; the accompanying coverage record distinguishes pages actually read from arguments independently checked.\par
\section*{Problem 1: resonance contradicts the decay in Lemma 4.8}
Classification: Major. Exact location: Lemma 4.8, p. 610, equation (4.38), and the cusp-end barrier paragraph in its proof. The necessary operator and homogeneous-solution hypotheses are on pp. 603, 605 and 607, equations (4.6), (4.8), (4.24)–(4.25) and Lemma 4.5.\par
\subsection*{The actual hypotheses and their leading coefficient}
Equation reference: eq:operator, eq:remainder. Set \(m=n+1\) and write functions in the \(\eta\) coordinate. Lemma 4.3 gives \[L_\infty=\frac1m D_n+\mathcal R,\qquad
 D_n=\eta^2\partial_{\eta\eta}-(n-1)\eta\partial_\eta-(n+1).
 \] The operator error is relative to the weighted derivatives, not an unweighted bounded perturbation. Definition 2.27, p. 570, and the explicit estimate (4.25), p. 607, give the bound \[|\mathcal R f|\leq C\eta^m
     \bigl(|f|+\eta|f'|+\eta^2|f''|\bigr)
 \] near the endpoint. Enlarging the right side to include \(|f|\) is harmless.\par
Equation reference: eq:indicial. Let \(U(\eta)=\widehat u(1-\eta)\), where \(\widehat u\) is exactly the function of Lemma 4.5. It solves \(L_\infty U=0\), is strictly positive on \((0,1)\), tends to \(1\) at the opposite endpoint, and satisfies \(U=O(\eta^m)\) at this endpoint. Uniform interior estimates for the rescaled ordinary differential equation on \([\eta,2\eta]\) also give \[|U|+\eta|U'|+\eta^2|U''|=O(\eta^m).\] Consequently \(D_nU=O(\eta^{2m})\). The indicial polynomial is \[D_n(\eta^r)=(r+1)(r-m)\eta^r,
 \] so the decaying and growing roots are \(m\) and \(-1\).\par
Equation reference: eq:quotient, eq:u-leading. For completeness, the coefficient of the decaying root cannot vanish here. Put \(U=\eta^m Z\). A direct calculation gives the useful exact identity \[D_n(\eta^m Z)
   =\eta^m\bigl(\eta^2Z''+(m+2)\eta Z'\bigr)
   =(\eta^{m+2}Z')'.
 \] Integrating \(D_nU=O(\eta^{2m})\) and using boundedness of \(Z\) removes the \(\eta^{-1}\) homogeneous branch. A second integration yields \[U(\eta)=A\eta^m+O(\eta^{2m})
 \] with the corresponding weighted first and second derivative estimates.\par
Here is an explicit uniqueness check for \(A\neq0\). If \(A=0\), write \(F=D_nU=-m\mathcal RU\). The two zero integration constants in [eq:quotient] (\#eq:quotient) give \[Z(\eta)=\int_0^\eta t^{-m-2}\int_0^t F(x)\,dx\,dt.\] For \[M(r)=\sup_{0<\eta\leq r}\eta^{-m}
              (|U|+\eta|U'|+\eta^2|U''|)<\infty,\] equation [eq:remainder] (\#eq:remainder) bounds \(F(x)\) by \(CM(r)x^{2m}\). The displayed integral and its first two weighted derivatives then imply \(M(r)\leq Cr^mM(r)\). For small \(r\) this forces \(U=0\); ordinary differential equation uniqueness extends this to the whole interval, contradicting \(\widehat u(0)=1\). Thus \(A\neq0\), and positivity gives \(A>0\). This rules out a vanishing leading coefficient as an escape from resonance.\par
\subsection*{The contradiction for the stated inhomogeneous solution}
Equation reference: eq:false-bound, eq:contradiction. Lemma 4.8 asks for \(V(\eta)=\widehat v(1-\eta)\) satisfying \[L_\infty V=U,\qquad
 |V|+\eta|V'|+\eta^2|V''|=O(\eta^m).
 \] Assume these exact requirements. Then [eq:remainder] (\#eq:remainder) gives \(\mathcal RV=O(\eta^{2m})\), and [eq:u-leading] (\#eq:u-leading) implies \[D_nV=mA\eta^m+O(\eta^{2m}).\] With \(V=\eta^m W\), identity [eq:quotient] (\#eq:quotient) becomes \[(\eta^{m+2}W')'=mA\eta^m+O(\eta^{2m}).\] Integrating gives \[W'=C\eta^{-m-2}+\frac{mA}{m+1}\eta^{-1}+O(\eta^{m-1}).\] Boundedness of \(W\) first forces \(C=0\), but a second integration then gives \[W(\eta)=\frac{mA}{m+1}\log\eta+C_0+O(\eta^m),
 \] which is unbounded because \(A>0\). Therefore no function satisfies [eq:false-bound] (\#eq:false-bound). This uses the actual forcing \(\widehat u\) and the paper’s actual operator error; it is not a counterexample obtained by replacing \(L_\infty\) with its model.\par
Equation reference: eq:log-term. The leading term required for a corrected decaying solution is \[V(\eta)=\frac{n+1}{n+2}A\eta^{n+1}\log\eta+O(\eta^{n+1}).
 \] Indeed \(D_n(\eta^{n+1}\log\eta)=(n+2)\eta^{n+1}\). In the paper’s dimension this is \[V(\eta)=\frac{3A}{4}\eta^3\log\eta+O(\eta^3).\] The logarithmic leading term is negative, so it does not contradict the desired sign \(\widehat v\leq0\). Adding a multiple of \(\widehat u\) changes the pure \(\eta^{n+1}\) term and cannot remove the logarithm. In particular, the adjustment imposing the Neumann condition at the other endpoint does not repair (4.38).\par
\subsection*{The failed barrier and the proof dependency}
Equation reference: eq:barrier. The barrier on p. 610 is of the form \[v_{\rm sub}=-C_1\eta^m+C_2\eta^{2m}.\] Its model image is exactly \[\frac1mD_nv_{\rm sub}=(2n+3)C_2\eta^{2m},
 \qquad L_\infty v_{\rm sub}=O(\eta^{2m}).
 \] It cannot dominate \(U\sim A\eta^m\) for all sufficiently small \(\eta\), for any fixed choices of the barrier constants. The proof’s asserted strictly positive constant lower bound omits this vanishing power. Making the endpoint interval smaller cannot change the order discrepancy.\par
The core dependency is explicit. Definition 4.9, p. 611, cuts off \(\widehat v\) and defines \(\widehat u_\sigma=L_\sigma\widehat v_\sigma\). The first sentence of p. 613 invokes (4.38) to prove the uniform bound \(|\widehat u_\sigma|\leq C\eta^3\) in Lemma 4.10. The dominated convergence argument on p. 614 uses that bound in Proposition 4.2. Theorem 4.12, p. 615, gives the uniform inverse modulo obstructions; p. 621 explicitly uses Proposition 4.2 to eliminate the nonzero limiting obstruction. The fixed-point construction then uses that inverse.\par
There are further direct uses in Section 5: the projected forcing estimate (5.13) and the gradient bound (5.18), p. 623, use (4.38); Claim 5.3’s projection estimate, p. 625, uses the same pure-power decay against the cusp volume density. Theorem 5.4, p. 626, removes the obstruction, and its estimate is used to deduce the Main Theorem on p. 628. Thus the defect affects a central analytic construction and its propagated estimates. It is classified as Major on this impact, irrespective of the eventual length of a repair.\par
A plausible repair begins by allowing \(\eta^{n+1}(1+|\log\eta|)\), with two weighted derivatives, and then rebuilding the cutoffs, uniform obstruction estimates, weighted pairings, projection bounds and nonlinear iteration. Logarithms are integrable in several of the displayed pairings and can sometimes be absorbed by a small power; these facts are evidence that a repair may be possible. They do not establish all uniform constants or preserve every stated exponent. Only the forced local logarithmic term is established here. The existence and sign properties with the left Neumann condition, the complete corrected gluing argument and the main theorem remain outside the verified repair. No counterexample to the Kähler–Einstein degeneration itself is claimed.\par
\section*{Problem 2: the expanded map has the wrong complex phase}
Classification: Minor. Exact location: Lemma 2.5, p. 550, the definition of \(\nu_\sigma\) and equations (2.18)–(2.20). The conjugates in Lemma 2.3 and (2.16) are essential.\par
Equation reference: eq:correct-map. Write \(F(z)=z_1^3+z_2^3+z_3^3\) and \(TY_\sigma=\{F=\sigma\}\). The base map is \[\Phi(w)_j=w_j+\nu(w)\overline w_j^{\,2}.\] Definition 2.4 correctly defines the scaled map by a composition. For a chosen cube root \(\lambda=\sigma^{-1/3}\) this gives \[\Phi_\sigma(z)_j
 =\lambda^{-1}\Phi(\lambda z)_j
 =z_j+\lambda^{-1}\overline\lambda^{\,2}
             \nu(\lambda z)\overline z_j^{\,2}.
 \] The lemma instead uses \(\lambda\nu(\lambda z)\) as the coefficient. In general \(\lambda^{-1}\overline\lambda^{\,2}\neq\lambda\). Their absolute values agree, but their phases need not.\par
\subsection*{An admissible complex-parameter check}
Equation reference: eq:a. Take \(\sigma=i\rho\) with \(\rho>0\) arbitrarily small, and \[\lambda=\rho^{-1/3}e^{-i\pi/6},\qquad
 z=(r,-r,0),\quad r>0.\] Then \(z\in TY_0\) and \(|\lambda z|>R\) for sufficiently small \(\rho\), so this point lies in the exact domain of Lemma 2.5. Choose the small positive root \(a\) of \[6r^2a-2a^3=\rho;
 \] it exists by the implicit function theorem, with \(a\sim\rho/(6r^2)\). Writing \(L=\rho^{-1/3}\) and \(q=a/(Lr^2)\), direct expansion gives \[F\bigl(\lambda z+q\,\overline{\lambda z}^{\,2}\bigr)
   =6qL^4r^4-2q^3L^6r^6=L^3\rho=1.\] Here the square is componentwise. The small root \(q\) is exactly the branch \(\nu(\lambda z)\) selected in Lemma 2.3.\par
Equation reference: eq:phase-witness. The composition [eq:correct-map] (\#eq:correct-map) therefore sends \(z\) to \((r+ia,-r+ia,0)\), whose cubic value is \(i\rho\). The printed formula instead sends it to \[(r+ae^{-i\pi/6},-r+ae^{-i\pi/6},0).\] The real part of this point’s cubic value is \[3\sqrt3\,r^2a>0.
 \] It is consequently not a point of \(TY_{i\rho}\). This countercheck works for arbitrarily small nonzero parameters in the stated complex family, and respects the large-radius domain of the base map. Restricting to positive real \(\sigma\) would hide the discrepancy, but is not the theorem’s parameter range.\par
\subsection*{Repair and boundary}
Replace the expanded coefficient by \[\nu_\sigma^{\rm corr}(z)
   =\lambda^{-1}\overline\lambda^{\,2}\nu(\lambda z).\] For real-coordinate derivatives, the chain rule and Lemma 2.3 give \[|\partial^k\nu_\sigma^{\rm corr}(z)|
 \leq C_k|\lambda|^{k+1}|\lambda z|^{-4-k}
 =C_k|\sigma|\,|z|^{-4-k}.\] Thus the stated decay estimates survive this correction. Changing the cube root by a third root of unity also leaves the corrected map unchanged: uniqueness of the small branch gives \(\nu(\zeta w)=\nu(w)\) and \(\zeta^{-1}\overline\zeta^{\,2}=1\). The defining composition (2.17) already contains the correct geometry. The error is therefore local and is classified as Minor: the expanded assertion fails to land in its target, while the composition and the norm estimates have an explicit repair. This does not certify every later coordinate computation in the gluing construction.\par
\section*{Independent consistency checks and their limits}
The family [eq:family] (\#eq:family) has affine critical points satisfying \(3z_j^2(2z_j^3+1)=0\). Their critical values are \(-k/4\) for \(k=0,1,2,3\). Hence for \(0<|\sigma|<1/4\) there is no affine singularity, and the central affine fiber has only the origin as a singular point. At infinity the leading Fermat sextic has no projective point at which all its coordinate derivatives vanish. The local map \(z_j\mapsto z_j(1+z_j^3)^{1/3}\) has derivative the identity at the origin and converts [eq:family] (\#eq:family) exactly to \(F=\sigma\). These elementary checks agree with the algebraic setup; they do not prove any metric convergence.\par
Equation reference: eq:horn. The radial equation (2.42), p. 555, with \(\psi',\psi''>0\) integrates to \[\frac{(\psi')^{n+1}}{n+1}=e^{\psi+a}+b,\] as in (2.79), p. 563. For \(b=-B<0\), putting \(y=t-T\) and \(u=\psi-\psi(T)\) gives \[\int_0^u(e^x-1)^{-1/(n+1)}\,dx
       =((n+1)B)^{1/(n+1)}y.\] Near \(y=0\) this yields the leading potential \[u\sim \frac{n^{(n+1)/n}}{n+1}B^{1/n}y^{(n+1)/n}.
 \] The total integral is finite and, after \(x\mapsto e^{-x}\), equals \[c(n)=\mathrm B\!\left(\frac1{n+1},\frac n{n+1}\right)
     =\frac{\pi}{\sin(\pi/(n+1))}.\] It gives the principal relation \(T\sim-(n+1)^{-1/(n+1)}c(n)B^{-1/(n+1)}\) in (2.89), p. 565. For \(n=2\), [eq:horn] (\#eq:horn) is \((2\sqrt2/3)B^{1/2}y^{3/2}\). Together with \(T=(2/3)\log|\sigma|\), p. 580, this explains the metric scale \(B^{1/2}\asymp|\log|\sigma||^{-3/2}\). This checks the principal scaling, not all normalization constants or remainder estimates in Section 2.\par
Figure 1, p. 544, and Figure 2, p. 564, are inspected for geometric orientation. Their schematic regions are not numerical evidence for an asymptotic estimate. The endpoint resonance is established from the equations and the weighted error bound, not inferred from those drawings.\par
\section*{Assessment and verification boundary}
Item | Assessment | Decisive check | Repair boundary\par
Lemma 4.8 decay | Major | Actual forcing gives a nonzero resonant logarithm | Uniform gluing estimates must be rebuilt\par
Lemma 2.5 phase | Minor | Admissible image misses \(TY_{i\rho}\) | Composition and norm estimates repaired\par
Algebraic setup | Consistent | Exact critical values and local coordinates | No metric conclusion\par
Principal neck scale | Consistent | First integral and horn expansion | Full errors and normalization unverified\par
Main Theorem | Unverified | Dependency analysis only | Independent acceptance pending\par
The original verification script was actually executed with exit code zero. It checked the general indicial polynomial, quotient identity, resonant coefficient and barrier order, the exact cubic phase witness, critical values and the leading horn coefficient. These symbolic identities corroborate the analytic arguments above; a finite test is not used as a substitute for an asymptotic proof. No author code was executed.\par
The detailed coverage record lists every source page read and every source image actually inspected. The full seven-region Ricci-potential estimates, all weighted Schauder and Liouville arguments away from the selected neck analysis, the complete fixed-point estimates after a logarithmic repair, and the no-other-bubble curvature argument are not independently certified. External existence and asymptotic theorems cited by the authors were not re-proved or imported as viewed-page evidence. No mathematical defect is inferred from an absence of code or supplements.\par
A bounded correction search and the publisher HTML corroborated identity and the relevant printed formulas; they do not prove that no correction exists. The staged manifest retains the acquisition searches, their limits and the exact source version. Actual parser output, the calculation script, exact TeX, compiler and dependency bytes, command and exit status, final PDF, page observations and the full support inventory are bound by the matching verification records. Scientific decision remains pending, the human-review list is empty, and this initial review is not independent acceptance.\par
9 X. Fu, H.-J. Hein and X. Jiang, A Continuous Cusp Closing Process for Negative Kähler–Einstein Metrics, Geom. Funct. Anal. 35 (2025), 542–632. Assigned publisher version of record, 91 pages. DOI: 10.1007/s00039-025-00708-y. Springer Nature, article identity and full HTML, https://link.springer.com/article/10.1007/s00039-025-00708-y, consulted for publication/version context and corroboration of the cited formulas.\par
\end{document}
